% ============================================================
\chapter*{Pr\'eface}
\addcontentsline{toc}{chapter}{Pr\'eface}
% ============================================================

\section*{Objectifs de ce cours}

Les bases de donn\'ees constituent l'un des piliers fondamentaux de l'informatique
moderne. Depuis les premiers syst\`emes de gestion de fichiers des ann\'ees 1960
jusqu'\`a l'explosion des bases NoSQL et des lacs de donn\'ees (\emph{data lakes}),
la capacit\'e \`a stocker, organiser et interroger des donn\'ees de mani\`ere fiable
et performante reste un enjeu central pour toute application informatique.

Ce cours s'adresse aux \'etudiants de deuxi\`eme ann\'ee de licence (L2) en
informatique, math\'ematiques-informatique ou sciences des donn\'ees. Il suppose
une familiarit\'e de base avec la programmation (Python) et les math\'ematiques
discr\`etes (ensembles, relations, logique propositionnelle).

\`A l'issue de ce cours, l'\'etudiant sera capable de~:
\begin{itemize}
    \item Concevoir un sch\'ema conceptuel (mod\`ele Entit\'e-Relation) et le
          traduire en sch\'ema relationnel normalis\'e.
    \item \'Ecrire des requ\^etes SQL complexes impliquant jointures, agr\'egations,
          sous-requ\^etes et fonctions fen\^etres.
    \item Comprendre et appliquer les formes normales (1NF \`a BCNF) pour
          \'eliminer les anomalies de mise \`a jour.
    \item Expliquer les propri\'et\'es ACID et les m\'ecanismes de contr\^ole
          de concurrence.
    \item Choisir et cr\'eer des index adapt\'es pour optimiser les performances.
    \item Concevoir un pipeline ETL simple avec Python et SQL.
    \item Comparer les mod\`eles NoSQL (document, cl\'e-valeur, colonne, graphe)
          au mod\`ele relationnel.
\end{itemize}

\section*{Organisation du cours}

Le cours est structur\'e en onze chapitres, progressant des fondements th\'eoriques
vers les applications pratiques~:

\begin{enumerate}
    \item \textbf{Introduction aux bases de donn\'ees relationnelles} ---
          Histoire, concepts fondamentaux, architecture d'un SGBD.
    \item \textbf{Mod\`ele Entit\'e-Relation} ---
          Conception de sch\'emas conceptuels, diagrammes E/R, passage au relationnel.
    \item \textbf{Alg\`ebre relationnelle} ---
          Op\'erateurs formels~: s\'election, projection, jointure, division.
    \item \textbf{SQL --- Requ\^etes de base} ---
          \texttt{SELECT}, \texttt{FROM}, \texttt{WHERE}, tri et filtrage.
    \item \textbf{SQL --- Jointures et sous-requ\^etes} ---
          Jointures internes, externes, crois\'ees ; sous-requ\^etes corr\'el\'ees.
    \item \textbf{SQL --- Agr\'egation et fonctions fen\^etres} ---
          \texttt{GROUP BY}, \texttt{HAVING}, \texttt{ROW\_NUMBER}, \texttt{RANK}.
    \item \textbf{Normalisation} ---
          D\'ependances fonctionnelles, 1NF \`a BCNF, d\'ecomposition sans perte.
    \item \textbf{Transactions, int\'egrit\'e, contraintes} ---
          Propri\'et\'es ACID, niveaux d'isolation, verrouillage.
    \item \textbf{Indexation et optimisation} ---
          B-arbres, index de hachage, plans d'ex\'ecution.
    \item \textbf{Pipelines de donn\'ees et ETL} ---
          Extraction, transformation, chargement ; int\'egration Python/SQL.
    \item \textbf{Bases NoSQL} ---
          Document, cl\'e-valeur, colonne, graphe ; th\'eor\`eme CAP.
\end{enumerate}

\section*{Approche p\'edagogique}

Chaque chapitre suit une structure coh\'erente~:

\begin{itemize}
    \item \textbf{Th\'eorie} --- D\'efinitions formelles, th\'eor\`emes et propositions
          avec d\'emonstrations ou justifications.
    \item \textbf{Exemples d\'etaill\'es} --- Chaque concept est illustr\'e par un
          ou plusieurs exemples concrets, souvent sur une base de donn\'ees
          \og Universit\'e \fg utilis\'ee comme fil rouge.
    \item \textbf{Code SQL ex\'ecutable} --- Les requ\^etes sont pr\'esent\'ees avec
          leur r\'esultat attendu, testables sur SQLite, PostgreSQL ou MySQL.
    \item \textbf{Int\'egration Python} --- Des scripts \texttt{sqlite3} et
          \texttt{pandas} montrent l'utilisation programmatique.
    \item \textbf{Exercices} --- Des exercices de difficult\'e croissante,
          allant de la v\'erification de compr\'ehension \`a la r\'esolution
          de probl\`emes ouverts.
\end{itemize}

\section*{Base de donn\'ees fil rouge}

Tout au long de ce cours, nous utilisons une base de donn\'ees
\textbf{Universit\'e} comportant les tables suivantes~:

\begin{codeblock}[title=\textbf{Sch\'ema de la base Universit\'e}]
\begin{minted}{sql}
CREATE TABLE Etudiants (
    id_etudiant  INTEGER PRIMARY KEY,
    nom          TEXT NOT NULL,
    prenom       TEXT NOT NULL,
    date_naissance DATE,
    filiere      TEXT
);

CREATE TABLE Cours (
    id_cours     INTEGER PRIMARY KEY,
    intitule     TEXT NOT NULL,
    credits      INTEGER DEFAULT 3,
    id_prof      INTEGER REFERENCES Professeurs(id_prof)
);

CREATE TABLE Professeurs (
    id_prof      INTEGER PRIMARY KEY,
    nom          TEXT NOT NULL,
    prenom       TEXT NOT NULL,
    departement  TEXT
);

CREATE TABLE Inscriptions (
    id_etudiant  INTEGER REFERENCES Etudiants(id_etudiant),
    id_cours     INTEGER REFERENCES Cours(id_cours),
    note         REAL,
    annee        INTEGER,
    PRIMARY KEY (id_etudiant, id_cours)
);
\end{minted}
\end{codeblock}

Ce sch\'ema sera enrichi au fil des chapitres pour illustrer des concepts
avanc\'es (h\'eritage, relations ternaires, historisation).

\section*{Donn\'ees d'exemple}

\begin{codeblock}[title=\textbf{Insertion des donn\'ees d'exemple}]
\begin{minted}{sql}
INSERT INTO Professeurs VALUES (1, 'Dupont', 'Marie', 'Informatique');
INSERT INTO Professeurs VALUES (2, 'Martin', 'Jean', 'Mathematiques');
INSERT INTO Professeurs VALUES (3, 'Durand', 'Sophie', 'Informatique');

INSERT INTO Cours VALUES (101, 'Bases de donnees', 4, 1);
INSERT INTO Cours VALUES (102, 'Algorithmes', 3, 2);
INSERT INTO Cours VALUES (103, 'Reseaux', 3, 3);
INSERT INTO Cours VALUES (104, 'Statistiques', 3, 2);

INSERT INTO Etudiants VALUES (1, 'Bernard', 'Alice', '2004-03-15', 'Info');
INSERT INTO Etudiants VALUES (2, 'Petit', 'Bob', '2003-11-22', 'Info');
INSERT INTO Etudiants VALUES (3, 'Moreau', 'Claire', '2004-07-08', 'Maths');
INSERT INTO Etudiants VALUES (4, 'Leroy', 'David', '2003-01-30', 'Info');
INSERT INTO Etudiants VALUES (5, 'Roux', 'Emma', '2004-09-12', 'Maths');

INSERT INTO Inscriptions VALUES (1, 101, 15.5, 2025);
INSERT INTO Inscriptions VALUES (1, 102, 13.0, 2025);
INSERT INTO Inscriptions VALUES (2, 101, 12.0, 2025);
INSERT INTO Inscriptions VALUES (2, 103, 16.0, 2025);
INSERT INTO Inscriptions VALUES (3, 102, 17.5, 2025);
INSERT INTO Inscriptions VALUES (3, 104, 14.0, 2025);
INSERT INTO Inscriptions VALUES (4, 101, 9.0, 2025);
INSERT INTO Inscriptions VALUES (5, 104, 18.0, 2025);
\end{minted}
\end{codeblock}

\section*{Outils recommand\'es}

\begin{itemize}
    \item \textbf{SQLite} --- L\'eger, sans serveur, id\'eal pour l'apprentissage.
          Disponible via \texttt{sqlite3} en ligne de commande ou le module
          Python \texttt{sqlite3}.
    \item \textbf{PostgreSQL} --- SGBD complet, conforme au standard SQL,
          recommand\'e pour les exercices avanc\'es (fonctions fen\^etres,
          transactions concurrentes).
    \item \textbf{DB Browser for SQLite} --- Interface graphique pour
          visualiser et manipuler des bases SQLite.
    \item \textbf{DBeaver} --- Client universel supportant PostgreSQL,
          MySQL, SQLite et d'autres.
    \item \textbf{Python 3.x} --- Avec les modules \texttt{sqlite3} (standard),
          \texttt{psycopg2} (PostgreSQL), \texttt{pandas} (analyse de donn\'ees).
\end{itemize}

\section*{Conventions typographiques}

\begin{itemize}
    \item Les mots-cl\'es SQL sont \'ecrits en \texttt{MAJUSCULES}~:
          \texttt{SELECT}, \texttt{WHERE}, \texttt{JOIN}.
    \item Les noms de tables et colonnes sont en \texttt{PascalCase} ou
          \texttt{snake\_case} selon le contexte.
    \item Les termes techniques apparaissant pour la premi\`ere fois sont
          en \emph{italique}.
    \item Les encadr\'es color\'es signalent~:
          \begin{itemize}
              \item les \textbf{d\'efinitions} et \textbf{th\'eor\`emes} importants,
              \item les \textbf{bonnes pratiques} \`a retenir,
              \item les \textbf{anti-patterns} \`a \'eviter,
              \item les \textbf{avertissements} sur les pi\`eges courants.
          \end{itemize}
\end{itemize}

\section*{Pr\'erequis}

\begin{itemize}
    \item \textbf{Programmation} --- Savoir \'ecrire un programme Python simple
          (variables, boucles, fonctions, fichiers).
    \item \textbf{Math\'ematiques} --- Notions d'ensembles, de relations,
          de produit cart\'esien, de logique propositionnelle
          ($\land$, $\lor$, $\neg$, $\Rightarrow$).
    \item \textbf{Syst\`eme} --- Savoir utiliser un terminal (ligne de commande).
\end{itemize}

\section*{Ressources compl\'ementaires}

\begin{itemize}
    \item A.~Silberschatz, H.~Korth, S.~Sudarshan, \emph{Database System Concepts},
          McGraw-Hill, 7\ieme{} \'edition.
    \item R.~Elmasri, S.~Navathe, \emph{Fundamentals of Database Systems},
          Pearson, 7\ieme{} \'edition.
    \item R.~Ramakrishnan, J.~Gehrke, \emph{Database Management Systems},
          McGraw-Hill, 3\ieme{} \'edition.
    \item Documentation officielle de PostgreSQL~: \url{https://www.postgresql.org/docs/}
    \item Documentation SQLite~: \url{https://www.sqlite.org/docs.html}
    \item W3Schools SQL Tutorial~: \url{https://www.w3schools.com/sql/}
\end{itemize}

\vfill
\begin{flushright}
\textit{Bonne lecture et bon apprentissage !}
\end{flushright}
